\documentclass[12pt]{article}
% Shared preamble and text for the data-request letters.
% Replace the bracketed fields before submitting any letter.

\usepackage[a4paper,top=2.4cm,bottom=2.4cm,left=2.8cm,right=2.4cm]{geometry}
\usepackage[utf8]{inputenc}
\usepackage[T1]{fontenc}
\usepackage{lmodern}
\usepackage{microtype}
\usepackage{setspace}
\usepackage{array}
\usepackage{booktabs}
\usepackage{enumitem}
\usepackage{tabularx}
\usepackage{xcolor}
\usepackage[hidelinks]{hyperref}

\setstretch{1.12}
\setlength{\parindent}{0pt}
\setlength{\parskip}{6pt}
\setlist[itemize]{leftmargin=1.5em,itemsep=2pt,topsep=2pt}
\setlist[enumerate]{leftmargin=1.7em,itemsep=2pt,topsep=2pt}
\newcolumntype{Y}{>{\raggedright\arraybackslash}X}
\newcolumntype{P}[1]{>{\raggedright\arraybackslash}p{#1}}

\newcommand{\NamaMahasiswa}{[NAMA LENGKAP MAHASISWA]}
\newcommand{\NIMMahasiswa}{[NIM]}
\newcommand{\NamaPembimbing}{[NAMA DAN GELAR DOSEN PEMBIMBING]}
\newcommand{\JudulSementara}{Diagnosis \textit{Borrowed Size} dan \textit{Agglomeration Shadow} antara Jakarta dan Pusat-pusat Sekunder di Kawasan Aglomerasi Jakarta}

\newcommand{\KopSurat}{%
  \begin{center}
    \fbox{%
      \begin{minipage}{0.88\textwidth}
        \centering
        \textbf{KOP SURAT RESMI FAKULTAS TEKNIK UNIVERSITAS GADJAH MADA}\\
        \small Ganti blok ini dengan kop resmi sebelum surat dikirim.
      \end{minipage}}
  \end{center}
  \vspace{4pt}
}

\newcommand{\IdentitasMahasiswa}{%
  \begin{tabularx}{\textwidth}{@{}P{0.30\textwidth}Y@{}}
    Nama & \NamaMahasiswa \\
    NIM & \NIMMahasiswa \\
    Program studi & Sarjana Perencanaan Wilayah dan Kota \\
    Departemen dan fakultas & Departemen Teknik Arsitektur dan Perencanaan, Fakultas Teknik, Universitas Gadjah Mada \\
    Dosen pembimbing & \NamaPembimbing
  \end{tabularx}
}

\newcommand{\PernyataanPenggunaanData}{%
Data yang dimohon akan digunakan untuk kepentingan pendidikan dan penelitian nonkomersial. Data tidak akan dijual, dialihkan, atau disebarluaskan dalam bentuk mentah. Hasil yang dipublikasikan akan berupa statistik agregat, peta turunan, atau keluaran lain yang diizinkan oleh pemilik data. Pengolahan dan penyimpanan data akan mengikuti ketentuan kerahasiaan, lisensi, pembatasan publikasi, dan perjanjian penggunaan data yang berlaku.
}

\newcommand{\PernyataanTanpaBiaya}{%
Apabila permintaan ini tidak dapat dipenuhi tanpa biaya, mohon kami diberi tahu terlebih dahulu. Kami tidak menyetujui pembuatan tabulasi atau pengadaan data berbayar tanpa persetujuan baru dari mahasiswa dan dosen pembimbing.
}

\newcommand{\AbstraksiPenelitian}{%
\textbf{Latar belakang.} Kegiatan ekonomi, perjalanan, layanan, dan perumahan di sekitar Jakarta berlangsung melampaui batas administrasi DKI Jakarta. Hubungan tersebut tidak otomatis menunjukkan bahwa pusat-pusat sekunder telah memperkuat fungsi lokalnya. Suatu wilayah dapat memperoleh akses ke pasar metropolitan, tetapi pada saat yang sama semakin bergantung pada Jakarta untuk pekerjaan, layanan, atau fungsi bernilai tinggi.

\textbf{Tujuan.} Penelitian ini mengkaji apakah integrasi dengan Jakarta berjalan bersama penguatan fungsi lokal, suatu pola yang konsisten dengan \textit{borrowed size}, atau bersama defisit fungsi dan ketergantungan yang asimetris, suatu pola yang konsisten dengan \textit{agglomeration shadow}. Penelitian tidak bertujuan mengusulkan perluasan wilayah DKI Jakarta, mengubah batas pemerintahan, atau membuktikan hubungan kausal hanya dari potret satu periode.

\textbf{Metode.} Kasus utama adalah sistem metropolitan Jabodetabek. Wilayah Statistik Metropolitan Indonesia digunakan sebagai sabuk pembanding. DKI Jakarta diperlakukan sebagai satu inti dengan batas administrasi tetap. Kecamatan menjadi unit utama jika resolusi data mendukung. Data pada tingkat kabupaten/kota, zona transportasi, raster, dan titik dipertahankan pada resolusi aslinya sebelum dilakukan agregasi yang dapat ditelusuri. Analisis membedakan data yang teramati, tabulasi agregat, hasil pemodelan, dan proksi.

\textbf{Cakupan wilayah dan waktu.} Baseline utama menggunakan tahun 2024. Data tahun lain digunakan sebagai konteks atau pemeriksaan perubahan, dengan tahun dan perbedaan unit ditampilkan secara terbuka. Wilayah mencakup DKI Jakarta, kabupaten/kota dalam kawasan utama Jabodetabek, pusat sekunder yang ditetapkan melalui prosedur yang terdokumentasi, serta sabuk pembanding yang diperlukan untuk menguji fungsi relatif.

\textbf{Jenis data dan variabel.} Penelitian memerlukan data penduduk, luas, kepadatan, integrasi dan perjalanan, fungsi layanan lokal, pekerjaan menurut lokasi tempat kerja dan sektor jika tersedia, aksesibilitas jaringan, serta karakteristik pusat sekunder. Podes digunakan untuk mengukur kehadiran dan akses layanan, bukan sebagai ukuran langsung produktivitas atau jumlah pekerjaan. Data komuter pada tingkat kabupaten/kota tidak akan dipaksa menjadi estimasi kecamatan. Data OD yang dimodelkan akan diberi label sintetis dan diuji terhadap kendala yang tersedia.

\textbf{Rencana hasil.} Hasil penelitian berupa profil integrasi dan fungsi lokal, tipologi kondisi relasional, peta gradien, uji sensitivitas bobot dan ambang, serta penjelajah skenario berbasis WebGL. Penjelajah tersebut menghitung ulang dan menampilkan hasil secara interaktif. Ia bukan sumber data waktu nyata, ramalan kebijakan, atau simulasi perubahan kewenangan pemerintahan.
}

\newcommand{\LampiranAbstraksi}{%
  \clearpage
  \section*{Lampiran 2. Abstraksi penggunaan data}
  \AbstraksiPenelitian
}

\newcommand{\TandaTanganDTAP}{%
  Hormat kami,\\[2.2cm]
  \parbox[t]{0.85\textwidth}{[NAMA PENANDATANGAN]}\\
  \parbox[t]{0.85\textwidth}{[JABATAN PENANDATANGAN SESUAI ALUR PERSURATAN DTAP]}
}

\newcommand{\TandaTanganBPS}{%
  Hormat kami,\\[2.2cm]
  \parbox[t]{0.85\textwidth}{[NAMA PENANDATANGAN]}\\
  \parbox[t]{0.85\textwidth}{[DEKAN FAKULTAS TEKNIK ATAU PEJABAT SETINGKAT YANG DITERIMA BPS]}
}


% Confirm the current records custodian before sending. If the receiving
% office does not control the records, ask it to forward the request or name
% the office that does.

\begin{document}
\KopSurat

\begin{flushright}
Yogyakarta, [TANGGAL SURAT]
\end{flushright}

\begin{tabularx}{\textwidth}{@{}P{0.20\textwidth}Y@{}}
  Nomor & [DIISI OLEH UNIT PERSURATAN] \\
  Lampiran & Dua lampiran \\
  Hal & Permohonan data agregat dan dokumentasi teknis JUTPI 2 dan JUTPI 3
\end{tabularx}

Yth. Pejabat Pengelola Informasi dan Dokumentasi\\
\mbox{[KEMENTERIAN/LEMBAGA PEMEGANG ARSIP JUTPI]}\\
c.q. [UNIT PENGELOLA ARSIP ATAU KONEKTIVITAS]\\
Jakarta

Dengan hormat,

Sehubungan dengan penyusunan Tugas Akhir pada Program Sarjana Perencanaan Wilayah dan Kota, Departemen Teknik Arsitektur dan Perencanaan, Fakultas Teknik, Universitas Gadjah Mada, kami mengajukan permohonan informasi untuk mahasiswa berikut.

\IdentitasMahasiswa

Penelitian berjudul sementara ``\JudulSementara''. Penelitian ini mengkaji hubungan antara DKI Jakarta dan pusat-pusat sekunder di kawasan Jabodetabek. Laporan JUTPI 2 dan JUTPI 3 menunjukkan bahwa survei transportasi, zona, jaringan model, dan pengukuran pergerakan pernah disusun untuk mendukung perencanaan transportasi kawasan tersebut. Penelitian kami memerlukan sebagian keluaran agregat dan dokumentasi teknisnya untuk memeriksa hasil yang dibangun dari sumber terbuka.

Permohonan ini dibatasi pada data antarzona yang telah diagregasikan, geometri zona, jaringan model yang dapat dibagikan, parameter perjalanan, serta tabel kalibrasi dan validasi. Kami tidak meminta rekaman telepon pintar, data GPS individu, ADS-MEILI, atau pengenal pribadi.

Jika Kementerian Koordinator tidak menguasai arsip yang dimohon, kami memohon agar permohonan ini diteruskan kepada unit atau instansi pemegang data, atau kami diberi petunjuk mengenai kanal yang tepat. Kami bersedia menerima bentuk agregat yang lebih kasar dan keluaran yang telah disaring oleh pemilik data.

\PernyataanPenggunaanData

\PernyataanTanpaBiaya

Demikian permohonan ini kami sampaikan. Atas perhatian dan bantuan Bapak/Ibu, kami mengucapkan terima kasih.

\vspace{8pt}
Tembusan:
\begin{enumerate}
  \item Dosen pembimbing;
  \item Ketua Program Studi Sarjana Perencanaan Wilayah dan Kota;
  \item Arsip.
\end{enumerate}

\TandaTanganDTAP

\clearpage
\section*{Lampiran 1. Rincian data yang dimohon}

\renewcommand{\arraystretch}{1.25}
\begin{tabularx}{\textwidth}{@{}P{0.12\textwidth}P{0.34\textwidth}Y@{}}
\toprule
Prioritas & Informasi atau data & Unit dan kelengkapan yang dimohon \\
\midrule
P0 & Matriks asal dan tujuan antarzona & Keluaran agregat JUTPI 2 dan/atau JUTPI 3, dengan tahun, zona, tujuan perjalanan, moda, periode, faktor ekspansi, dan keterangan apakah data teramati atau termodelkan. \\
P0 & Geometri dan crosswalk zona & Kode dan nama 349 TAZ jika tersedia, geometri, sistem koordinat, tahun batas, serta padanan ke desa atau kelurahan, kecamatan, dan kabupaten atau kota. \\
P0 & Jaringan model dan parameter perjalanan & Ruas, simpul, moda, waktu atau biaya perjalanan, fungsi hambatan, tahun jaringan, dan aturan konektivitas yang boleh dibagikan. \\
P1 & Kalibrasi dan validasi & Distribusi panjang perjalanan, kecocokan terhadap nilai marginal, galat, sampel atau survei yang digunakan, dan keterbatasan model. \\
P1 & Tabulasi pusat sekunder & Arus dari dan menuju DKI Jakarta serta tiga sampai lima pusat sekunder prioritas, tanpa pengenal individu dan tanpa jejak lokasi mentah. \\
\bottomrule
\end{tabularx}

\vspace{8pt}
Jika data utama tidak dapat diberikan, kami tetap menerima geometri dan kamus zona, dokumentasi metode, parameter, atau tabel validasi. Data tersebut akan digunakan untuk mengaudit OD sintetis, bukan untuk menyatakan bahwa hasil model merupakan arus teramati.

\begin{itemize}
  \item Kami tidak meminta GPS personal, ADS-MEILI, rekaman telepon pintar, atau data individu.
  \item Kami tidak meminta data yang harus dibuka melanggar perjanjian kerahasiaan atau hak pihak ketiga.
  \item Kami bersedia menerima hasil agregat yang sudah disaring dan ketentuan publikasi turunannya.
\end{itemize}

\LampiranAbstraksi

\end{document}
